\documentclass[12pt]{article}

\usepackage{fancyhdr}
\usepackage{geometry}
\usepackage{ucs}
\usepackage[utf8x]{inputenc}
\usepackage[T1]{fontenc}
\usepackage[ngerman]{babel}
\usepackage{amsmath,amssymb,amstext}
\usepackage{hyperref}
\usepackage{cancel}
\usepackage{dsfont}
\usepackage{physics}
\usepackage{lmodern}
\usepackage{enumerate}
\usepackage{enumitem}
\usepackage{graphicx}
\usepackage{listings, color}
\usepackage[labelfont=bf]{caption}
\usepackage{titling}


\lstset{basicstyle=\scriptsize,breaklines=true} %Quellcode mit Umlauten und ganz klein

\lstset{literate=
  {Ö}{{\"O}}1
  {Ä}{{\"A}}1
  {Ü}{{\"U}}1
  {ß}{{\ss}}2
  {ü}{{\"u}}1
  {ä}{{\"a}}1
  {ö}{{\"o}}1
}


%Geometrie----------------------------------------------------------------------------------------------------------

\geometry{a4paper, top=25mm, left=15mm, right=15mm, bottom=25mm,headsep=10mm, footskip=10mm}
\pagestyle{fancy}
\setlength{\parindent}{0pt} %Zeileneinrückung

\fancyhf{} %Setzt voreingestellte Kopf-und Fußzeilen-Eigenschaften zurück

\lhead{\nouppercase{\leftmark}}
\chead{}
\rhead{\thepage}

\lfoot{}
\cfoot{}
\rfoot{}

\title{\vspace{0cm}{\Huge Fortgeschrittenen-Praktikum I:\\ \vspace{1cm}  Halbleiter}}
\author{Saskia Bondza\\Simon Stephan}
\date{Durchgeführt am 07.09.2016}

\pretitle{%
  \begin{center}
  \LARGE
  \includegraphics[width=6cm,]{figures/siegel}\\[\bigskipamount]
}
\posttitle{\end{center}}

%neue Commands----------------------------------------------------------------------------------------------------------
\newcommand{\nab}{\vec{\nabla}} %direkter Befehl mit Vektorpfeil
\newcommand{\gra}[3][0.7]{
	\begin{minipage}[h!]{\textwidth}
		\centering
		\includegraphics[width=#1\textwidth]{figures/#2.png}
		\captionof{figure}{#3}
	\end{minipage}
	\vskip 30 pt
	}


%Titel,Inhalt----------------------------------------------------------------------------------------------------------

\begin{document}
\pagenumbering{gobble} %verstecke Seitenzahl
\maketitle
\newpage

\thispagestyle{empty}
\tableofcontents
\newpage

%Schreiben----------------------------------------------------------------------------------------------------------
\pagenumbering{arabic} %verstecke Seitenzahl
\section{Einleitung}

In diesem Experiment soll in drei verschiedenen Versuchsteilen die Funktionsweise von Halbleitern und Halbleiterdetektoren untersucht werden. Silizium wird hier als Beispiel für einen reinen Halbleiter untersucht und Kadmiumtellurid als Beispiel für einen Verbindungshalbleiter.
In \textbf{Teil I} des Versuchs wird die Bandlückenenergie, wichtiges Charakteristikum von Halbleitern, gemessen, die die Auflösung eines Halbleiterdetektors bestimmt.\\
In \textbf{Teil II} des Versuchs geht es um die Beobachtung der Ausbreitung
von Ladungsträgern im Halbleiter mit dem Haynes-Shockley-Experiment und in \textbf{Teil III} des Experiments werden mit Halbleiterdetektoren radioaktive Spektren aufgenommen und untersucht, wie gut sich Halbleiterdetektoren hierfür eignen.


\newpage
\section[Theoretische Grundlagen]{Theoretische Grundlagen}

Dieser Teil orientiert sich an [].



\subsection{Interferenz von Photonen, optisches Gitter}

Trifft eine ebene Welle auf einen Spalt, so ist nach dem Huygene'schen Prinzip jeder Punkt der Wellenfront Ausgangspunkt einer Elementarwelle, wobei sich die neue Wellenfront durch Überlagerung sämtlicher Elementarwellen ergibt. Trifft eine ebene Welle auf mehrere Spalte, ein optisches Gitter, beobachtet man hinter dem Gitter Interferenz, ein charakteristisches Wellenphänomen.  Ist der Gangunterschied zweier Wellen ein ganzzahliges Vielfaches ihrer Wellenlänge kommt es zur konstruktiven Interferenz, ist der Gangunterschied ein ungerades, halbzahliges Vielfaches der Wellenlänge kommt es zur destruktiven Interferenz. Es gilt also folgender Zusammenhang für die konstruktive Interferenz am Gitter:

\begin{align*}
& n\cdot\lambda=d\cdot\sin(\theta)
\end{align*}
Hierbei ist $n$ die Beugungsordnung, $\theta$ der Einfallswinkel, $\lambda$ die Wellenlänge und $d$ der Spaltabstand.
Man unterscheidet zwei Arten von optischen Gittern: Transmissionsgitter und Reflexionsgitter; Transmissionsgitter bestehen aus parallelen Spalten, Reflexionsgitter aus reflektierenden Graten. Hier handelt es sich um ein Reflexionsgiterr. Wir betrachten nun nur die erste Beugungsordnung. Für verschiedene Anstellwinkel des optischen Gitters ergibt sich dann die Energien der Photonen zu:

\[E(\varphi)=\frac{hc}{2d\sin(\varphi)cos(\frac{\psi}{2})}\]

In 4.1 wird genauer erklärt wie sich ein optisches Gitter so als Spektrometer verwenden lässt.

\subsection{Halbleiter}

\subsection{Das Bändermodell}

Bei einem einzelnen Atome sind die Energieniveaus der Elektronen diskret. Nähert man nun zwei Atome so nah an, dass sich ihre Orbitale überlappen, müssen sich die Energieniveaus gemäß dem Pauliprinzip aufspalten. In einem Kristall wechseln nun der Art viele Atome miteinander, dass auf diese Art und Weise Energiebänder enstehen (s. Abbildung ). Das energetisch höchste Band, das im Grundzustand besetztist heißt Valenzband, das erste unbesetzte Band heißt Leitungsband. Der Energieabstand zwischen diesen beiden Bändern wird als Bandlücke bezeichnet. Bei $T = 0$K ist das Valenzband voll besetzt während im Leitungsband keine Ladungsträger sind. Durch Anlegen einer Spannung oder Absorption von Photonen können Elektronen jedoch in das Leiterband gelangen. 
Die charakteristische Größe der Bandlückenenergie definiert ob ein Stoff ein Isolator, Halbleiter oder Leiter ist. Isolatoren haben eine besonders große Bandlüclenenergie, die mehrere $eV$ beträgt. Halbleiter haben typischerweise eine Bandlücke von weniger als einem $eV$, sie isolieren nur bei sehr niedrigen Temperaturen. Bereits thermische Anregungen reichen bei höheren Temperaturen für die Elektronen aus, um ins Leitungsband gehoben zu werden, Zurück bleibt bei diesem Vorgang ein positiv geladenes Loch, man spricht von einem Elektron-Loch-Paar. Bei Leitern überlappen sich beide Energiebänder, sodass stets genug Ladungsträger im Leitungsband sind, um einen Stromfluss zu gewährleisten.

\subsection{Bewegung von Ladungsträgern im Halbleiter}

\subsection{Direkte und indirekte Halbleiter}

Untersucht man die Energien der Bänder in Abhängikeit des Impulses ergeben sich komplizierte Verläufe, aus denen sich eine natürliche Unterscheidung von Halbleitern ableiten lässt: Liegt das Maximum des Valenzbandes beim gleichen Impuls wie das Minimum des Leitungsbandes spricht man von direkten Halbleitern, liegen Maximum und Minimum bei verschiedenen Impulsen, spricht man von indirekten Halbleitern.
Bei direkten Halbleitern können Elektronen einfach durch die Aufnahme der Bandlückenenergie (Übertrag durch ein Photon von ausreichender Energie $E_{Photon}=\hbar\omega\geq E_g$) ins Leitungsband gehoben werden während bei indirekten Halbleitern durch die Verschiebung der Extrema um $\Delta p$ ein Elektron zusätzlich zur Energieaufnahme auch noch seinen Impuls ändern muss. Dies geschieht durch Erzeugung bzw. Vernichtung von Phononen, d.h. Gitterschwingungen.

\subsection{Extrinsische Halbleiter und Dotierung}

In der Theorie wird meist von intrinsischen Halbleitern ausgegangen, also perfekten Kristallen ausgegangen.  Diese lassen sich allerdings technisch nicht erzeugen. Kristalldefekte, z.B.Beschädigte Elementarzellen, Verschiebung ganzer Kristallebenen oder die Verunreinigung durch Fremdatom beeinflussen die Kristallstruktur und können zur Bildung von Energieniveaus innerhalb der Bandlücken führen, was einen nicht zu vernachlässigenden Effekt auf die Elektronen-Loch-Rekombination hat, und somit die Leitfähigkeit beeinflusst.\\
Normalerweise geht ein Halbleiteratom Bindungen mit vier Nachbaratomen ein. Bei der Verunreinigung mit Fremdatomen z.B. Phosphor mit fünf Valenzelektronen oder Aluminium mit drei Valenzelektronen steht ein quasi-freier Ladungsträger zur Verfügung: Bei Phosphor ist dies das fünfte, schwachgebundene Elektron, bei Aluminium steht ein zusätzliches Loch als Ladungsträger zur Verfügung. Bei Fremdatome, die ein zusätzliches Elektron beisteuern, wie in diesem Beispiel Phosphor, spricht man von Donatoren. Bei solchen, die ein zusätzliches Loch beisteuern, wie hier Aluminium, spricht man von Akzeptoren.
Oft wird diese Verunreinigung absichtlich vorgenommen (typischerweise ca. 1 Fremdatom auf $10^6$Atome), da sie sich positiv auf die Leitfähigkeit des Halbleiters auswirken kann. Diesen Vorgang nennt man Dotierung. Donatoren-Halbleiter nennt man n-Typ, Akzeptor-Halbleiter heißen p-Typ.
\subsection{Halbleiter-Diode}

Bringt man eine Schicht n-Typ-Halbleiter und eine Schicht p-Typ-Halbleiter zusammen, spricht man von einer Diode. Die quasie-freien Ladungsträger der n-Typ-Schicht, d.h. die Elektronen, diffundieren in die freien Löcher der p-Typ-Schicht. Im Übergangsbereich löschen sich die freien Ladungsträger also durch diesen Vorgang aus. 
Nach diesem Diffusionsprozess bleiben in der n-Typ-Schicht positiv geladene Ionen und in der p-Typ-Schicht negativ geladene Ionen zurück (Verarmungszone), die sich nicht frei bewegen können, da sie in die Gitterstruktur des Kristalls fest eingebunden sind. Es ensteht so ein elektrisches Feld (Kontakt-Potential mit Spannung $U_{bi}$). 

Durch Anlegen einer äußeren Spannung bei der der positive Pol an die n-Schicht angeschlossen wird, und der negative Pol an die p-Schicht vergrößert sich die Verarmungszone und die Ladungsträger müssen ein größeres Spannungsgefälle überwinden: die Diode ist in Sperrrichtung gepolt. Bei Anlegung einer äußeren Spannung mit umgekehrter Polung können die überschüssigen Elektronen in der p-Schicht problemlos zum positiven Pol wandern, die Diode ist dann in Durchlassrichtung gepolt.
Dies bildet die Grundlage für Halbleiterdetektoren. Die Grundidee von Halbleiterdetektoren ist gerade eine Halbleiterdiode in einen Stromkreis einzubinden.
\subsection{Wechselwirkung von Strahlung mit Materie}

Die Wechselwirkung von Strahlung mit Materie wird im Wesentlichen durch drei verschiedene Effekte beschrieben: \textbf{Photoeffekt}, \textbf{Comptoneffekt} und \textbf{Paarbildung}.
Der Effekt der Paarbildung tritt für Energien über $1,022$ MeV auf, die in diesen Versuchen nicht erreicht werden, sodass wir diesen Effekt hier nicht weier erläutern.

\paragraph{Photoeffekt} Man unterscheidet beim Photoeffekt zwischen innerem und äußerem Photoeffekt.
Beim Photoeffekt dringt ein Photon in das Atom ein und überträgt seine gesamte Energie an ein Elektron der inneren Schalen. Dabei wird Energie auf dieses Elektron übertragen. Ist die Energie des Photons größer als die Bindungsenergie des Elektrons wird es aus der Atomhülle befreit und erhält kinetische Energie. Die hier entstandene Lücke wird über Abstrahlung eines $\gamma$-Quants oder eines Auger-Elektrons wieder gefüllt. Man spricht vom äußeren Photoeffekt.
Beim inneren Photoeffekt treffen Photonen mit einer Energie  $\geq E_g$ auf ein Elektron, und heben dieses vom Valenzband in das Leiterband. Hier durch wird die Leitfäjigkeit des Halbleiters erhöht.
\paragraph{Comptoneffekt} Beim Compton-Effekt trifft ein einfallendes Photon auf ein freies oder nur leicht gebundenes Elektron und überträgt einen Teil seiner Energie bzw. seines Impulses auf dieses und verliert selber an Energie in dem die Wellenlänge erhöht wird. Die Stärke dieses Effektes hängt vom Emissionswinkel ab, gestreute Photonen können weiter mit anderen Elektronen wechselwirken, beispielsweise in Form des inneren Photoeffekts. Dieser Prozess findet meist bei Energien zwischen $200$ keV und $5$ MeV statt.

Trifft hochenergetische Strahlung, z.B. in Form eines Lasers, auf einen Halbleiter, ensteht durch die oben beschriebenen Effekte eine Ladungsträgerwolke.

In Materie fällt die Intensität $I_d$ der $\gamma$-Strahlung exponentiell mit der Strecke $d$ ab:
\[I_d=e^{\mu d}\]
Dabei ist $\mu$ der Absorptionskoeffizent des Materials und $I_0$ die Anfangsintensität.

\subsection{Theorie zum Haynes-Shockley Experiment}

Im zweiten Veruchsteil des Experiments wird ein Laserpuls auf eine Germaniumprobe geschossen, wodurch eine Ladungswolke ensteht.Um das Verhalten dieser Ladungsträgerwolke zu beschreiben stellen wir folgende Überlegungen an: Aus der Kontinuitätsgleichung $\vec{\nabla}\vec{j}+\dot{n}=0$ und zusätzlichen Rekombinationen (siehe \cite{staat}), ergibt sich:
\begin{align*}
&\frac{\partial n}{\partial t}=-\vec{\nabla}\vec{j_n}\frac{n-n_o}{\tau_n}=\vec{E}\mu_n\vec{\nabla}n+D_n\triangle n-\frac{n-n_o}{\tau_n}&\mbox{(analog für die Löcher)}
\end{align*}
wobei  $\tau_n$ die  Rekombinationszeit ist und $n_0,p_0$ sind die Grundkonzentrationen, für die gilt: $n(\vec{x})=n_0+\bar{n}(\vec{x})$

Diese Differentialgleichung ist gekoppelt. Wir vereinfachen sie daher durch Einfhren allgemeiner Größen zu einer entkoppelten DGL:
\[\frac{\partial c}{\partial t}=\mu\vec{E}\cdot\nab c+D\triangle c-\frac{c-c_0}{\tau}\]
wobei :
\begin{itemize}
	\item $c=$  ambipolare Ladungsträgerdichte.
	\item $c_0=$  Grundkonzentration von $c$.
	\item $D=$  Diffusionskonstante.
	\item $\mu=$ Beweglichkeit.
\end{itemize}
Man kann zeigen, dass folgende Funktion diese Differenzialgleichung löst:
\[c(x,t)=Ce^{-\frac{t}{\tau_n}}\frac{1}{\sqrt{4\pi D_n t}}e^{-\frac{(x-\mu_n E t)^2}{4D_n t}}\]
Bei genauerem Hinsehen erkennt man schnell, dass es sich im Wesentlichen um eine Gaußfunktion handelt. Es macht daher Sinn, für die Auswertung dieses Versuchsteils Fit-Funktionen folgender Form zu verwenden:
\[c(x,t)=A(t)\cdot\frac{1}{\sqrt{2\pi\sigma^2}}\cdot e^{-\frac{(x-x_c(t))^2}{2\sigma^2}}\]

mit:
\begin{itemize}
	\item $A(t)=Ce^{-\frac{t}{\tau_n}}$
	\item $x_c(t)=\mu_n E t$
	\item $\sigma(t)=\sqrt{2 D_n t}$
\end{itemize}

 
\newpage
\section[Versuchsaufbau- und Durchführung]{Versuchsaufbau- und Durchführung} 


\subsection{Teil I}


\gra{TeilI_Aufbau}{Versuchsaufbau Teil I: Messen der Bandlückenenergie}

Abbildung [] zeigt den schematischen Versuchsaufbau des Versuchs. Am Anfang des Strahlengangs steht eine Lampe als Strahlungsquelle, die weißes Licht

\subsection{Teil II}

\subsubsection{Versuchsaufbau}

\gra{TeilII_Aufbau}{Versuchsaufbau Teil II: Bewegung der Ladungsträgerwolke}

In Abbildung [] ist der Versuchsaufbau zum Haynes-Shockley-Versuch zu sehen. Licht eines gepulstes Lasers wird hier durch ein Glasfaserkabel auf eine Germaniumprobe ($40\times 40\times 300\mathrm{mm}^3$, p-dotiert) geleitet.  Hierdurch ensteht eine Ladungsträgerwolke, welche durch die angelegte Spannung $U$ durch den Halbleiter wandert. Wir beobachten diese Wolke mit Hilfe eines Oszilloskops und einer Sonde.
Um ordentlich messen zu können, ist es erforderlich, dass die Frequenz des gepulsten Lasers und die des Oszilloskops synchronisiert sind, was über die Verbindung des Trigger-Eingangs des Oszilloskops mit dem Laser bezweckt wird.

Da die Probe wieder abkühlen muss, ist die Spannung $U$ nicht konstant. Viel mehr werden $98\%$ der Zeit einer Periode für das Abkühlen der Probe verwendet während der eigentliche Prozess in den restlichen $2\%$ der Zeit stattfindet.
Da die Treiberspannung deutlich größer ist als die Ladungsträgerspannung, muss diese mit einem \glqq shifted output\grqq\ herausgefiltert werden.

\subsubsection{Versuchsdurchführung}

In diesem Teil werden zwei verschiedene Messreihen aufgenommen. In einer ersten Messreihe wird der Abstand zwischen Lichtkabel und Nadel variiert, während die Spannung konstant bleibt. Es werden zehn Messungen für Abstände zwischen ca. $2-10$ mm durchgeführt. In einer zweiten Messreihe wird nun bei einem konstanten Abstand von ca. $4$ mm (inklusive Offset) die angelegte Spannung variiert. Hier werden elf Messungen für Spannungen zwischen $20$ V und $50$ V vorgenommen. Beide Messreihen sollten Messpunkte in möglichst regelmäßigen Abständen enthalten.
\subsection{Teil III} 

\subsubsection{Aufbau}
\gra{TeilIII_Aufbau}{Versuchsaufbau Teil III: Halbleiterdetektoren}

Kernstück dieses Versuchsaufbaus ist der Halbleiterdetektor. Es stehen für diesen Versuchsteil zwei verscheidene Halbleiterdetektoren zur Verfügung:

\begin{itemize}
	\item Eine $n-n^+$-Silizium-Diode
	\item Ein CdTe-Kristall
\end{itemize} 

Es werden mit beiden Halbleiterdetektoren je die Spektren von  $^{57}$Co und $^{241}$Am aufgenommen. Durch Einstellen einer festen Messdauer wird eine einheitliche Normierung auf eine feste Messdauer überflüssig. \\
Trifft ein durch den Zerfall enstehendes $\gamma$-Quant auf den Halbleiter, so kann es seine Energie durch den Photoeffekt bzw. den Comptoneffekt ganz oder teilweise auf ein Elektron übertragen, welches so vom Valenzband in das Leiterband gehoben werden kann. Dieses kann durch Stöße seine Energie an weitere Elektronen übertragen. Die Anzahl der zusätzlichen Ladungsträger ist dabei proportional zur Ursprungsenergie des Photons.
Durch Anlegen einer Spannung kann so ein Strompuls gemessen werden. Dieser Strompuls wird vom Vorverstärker in einen Spannungssignal umgewandelt. Dieses Spannungssignal führt dann zum Aufladen eines Kondensator im Shaping Amplifier. Nachdem sich dieser wiederum über einen Widerstand entladen hat, wird das Signal über zwei RC-Filter gedämpft. Man erhält so einen scharfen Spannungspeak dessen Amplitude weiter proportional zur Ursprungsenergie des Photons ist. Der Multi Channel Analyzer ordnet damit verschiedene Energien den verschiedenen Channels zu und leitet dies an den Computer weiter mit dem wir nun schließlich Zerfallspektren aufzeichnen können.

\subsubsection{Durchführung}

Der Halbleiterdetektor wird in den vorgesehenen Kasten eingebaut und die radioaktive Probe daraufgestellt. Mit dem Computer-Programm ADMCA werden die Zerfalls-Spektren aufgenommen. Es werden vier Messungen durchgeführt mit einer Messzeit von jeweils einer Stunde:

\begin{itemize}
	\item Spektrum von $^{57}$Co mit Si-Detektor
	\item Spektrum von $^{57}$Co  CdTe-Detektor
	\item Spektrum von $^{241}$Am  mit Si-Detektor
	\item Spektrum von $^{241}$Am mit CdTe-Detektor
\end{itemize} 


\newpage
\section{Auswertung}



\newpage
\section{Zusammenfassung/Diskussion}




\subsection{Diskussion}\label{Diskussion}


\newpage
\section{Anhang}






%\subsubsection{$\alpha$-Plateau Samarium}
%\lstinputlisting{data/Americium_1.TKA}


%\newpage
%\subsection{Quellcode (MATLAB)}
%\lstinputlisting[language=MATLAB]{Rohdaten/alpha.m}


%\begin{minipage}{\textwidth}
%\centering
%\includegraphics[width=0.9\textwidth]{figures/IMG_20151002_141014.jpg}
%\end{minipage}

\newpage
\listoffigures

%Literatur----------------------------------------------------------------------------------------------------------

%\cite{les}
\newpage
\thispagestyle{empty}
\begin{thebibliography}{9}

\bibitem{anleitung}
 (Zerfall von $^{57}$Co): http://hacol13.physik.uni-freiburg.de/fp/Versuche/FP1/FP1-6-KurzeHalbwertzeiten/Anleitung.pdf
  
\bibitem{physchem}
	Gerd Weldler, Hans-Joachim Freund,
	\emph{Lehrbuch der Physkalischen Chemie, Band 1},
	6th edition,
	John Wiley \& Sons,
	2012

\bibitem{nuclear}
Bernard L. Cohen,
\emph{Concepts of Nuclear Physics},
 New York usw.: McGraw-Hill 1971
  
%\bibitem{molmasse}
%  \emph{http://www.convertunits.com/molarmass/<ELEMENTNAME AUF ENGLISCH>}, Stand 28.09.2015
  

\end{thebibliography}

\end{document}